\section{已解决到哪一步：逐项对应原始论文}
\label{sec:guide-solved}
以下每项给出作者引文、年份和稳定文献键。点击引文可进入参考文献，再由“全文核读档案”进入第7章；用文献键搜索HTML索引可直接打开本地原文及图号记录。
这里的“已解决”指在明确条件下建立了可检验机制或获得了可靠结果，不表示所有几何和参数下均已解决。

\begin{longtable}{P{29mm}P{63mm}P{56mm}}
\toprule 问题 & 结果与具体论文 & 证据边界\\\midrule\endfirsthead
\toprule 问题 & 结果与具体论文 & 证据边界\\\midrule\endhead
为何必须区分离子与电子流 &
\citet{Ren2005}（2005，Ren2005）图2、4测到MRX四极磁场及碰撞性关联，支持解耦电流系统；\citet{Birn2001}（2001，Birn2001）提供相关跨模型基准。 &
Ren2005的电子流及Hall电场含反演/估计，未闭合X点欧姆预算；不能写成Hall项直接断线已获验证。\\
已启动重联能否得到快演化 &
\citet{Birn2001}（2001，Birn2001）图1中Hall、混合与PIC通量曲线接近，建立了有限种子二维非线性阶段的基准。 &
强种子刻意跳过线性触发；电子层和喷流宽度仍不同。该文不证明自然启动或普适0.1。\\
同一参数下为何可有快慢两支 &
\citet{Cassak2005}（2005，Cassak2005）图3、4给出不同历史及迟滞；\citet{Cassak2010}（2010，Cassak2010）进一步在关闭电子惯性的Hall算例中保留双稳态。 &
依赖具体片层、边界和耗散；2010现象学模型未自洽选定层长，不能给出全部环境的触发阈值。\\
MAST双峰密度怎样形成 &
\citet{Stanier2013}（2013，Stanier2013）将Hall密度响应、环形内侧压缩和实验诊断弦结合，解释双峰及外峰随合并消失的顺序。 &
是形态和顺序的解释；峰值、时钟未逐时定量匹配，且采用固定温比和增强耗散。\\
X线为何单向或双向扩展 &
\citet{Shepherd2012}（2012，Shepherd2012）图3--6显示弱导场主要向电子流动方向扩展，强导场转为双向近导场Alfvén传播。 &
只检验了电子承担初始全部电流等条件；阈值和速度有诊断误差，后期斜模破坏简单图像。\\
流体能否改善电子层描述 &
\citet{Wang2015}（2015，Wang2015）图4、8--10显示十矩可改善部分局地电子流和压力结构；\citet{F611}（2016，F611）图1--7用各向异性流体电子再现长电流层。 &
Wang2015未重现完整时史及所有压力峰值；F611是混合动理学离子模型。这些不是标准标量Hall已完全解决电子层的证据。\\
\bottomrule
\end{longtable}

例如第一行的实验验证、第二行的基准一致和最后一行的闭合改善，属于不同类型的进展。不能把它们叠加成“某一套Hall代码已经同时完成全部验证”。

\section{未解决什么：原文缺口与本报告判断分别标注}
\label{sec:guide-unsolved}
以下问题均对应已核读论文。本节所说“未解决”限于这些论文提供的证据，不能据此断言所有后续文献都没有答案。“原文缺口”指文中明确保留或结果直接显示的不足；“本报告综合判断”是跨论文比较后提出的研究问题，不冒充作者原话。

\subsection{电子扩散区的绝对层宽、速度和欧姆预算}
\textbf{对应文章：}\citet{L119}（2008，L119，
\emph{Experimental study of the Hall effect and electron diffusion region during magnetic reconnection in a laboratory plasma}）；
\citet{Wang2015}（2015，Wang2015）。

\textbf{原文缺口：}L119图13的电场预算没有测全压力张量与涨落项；图16--18的电子层较所比较二维模型更宽、排出更慢，经验摩擦估计仍有差距。
Wang2015图4的重联时史不同；图8--10在等磁通、不同时间比较时，部分非对角压力峰值及下游欧姆项仍不匹配。
\textbf{本报告综合判断：}同为电子尺度的趋势一致，尚未构成对层厚、速度和全部非理想项的同时验证；下一步应在相同诊断定义和仪器平均下比较完整预算。

\subsection{同一压力/热流闭合能否跨算例迁移}
\textbf{对应文章：}\citet{Ng2020}（2020，Ng2020，
\emph{An improved ten-moment closure for reconnection and instabilities}）；
\citet{F611}（2016，F611）；
\citet{Closure2026}（2026，Closure2026）。

\textbf{原文缺口：}Ng2020图1--7及附录A显示，梯度闭合可改善岛合并和部分不稳定性，电子各向异性及跨场热流仍有误差；F611的Harris实现需要二维磁通标签及参考密度修正；Closure2026恢复了部分三维不稳定性，却仍有增长时序、各向异性位置与范围的差异。
\textbf{本报告综合判断：}尚未建立一套无需逐算例调整、又能同时预测率、层结构、热流和三维模态的通用闭合。由某一算例改善不能推出跨几何适用。

\subsection{局地快率能否预测全局完成时间和尺寸效应}
\textbf{对应文章：}\citet{Birn2001}（2001，Birn2001）；
\citet{L241}（2017，L241，
\emph{The role of guide field in magnetic reconnection driven by island coalescence}）；
\citet{Stanier2013}（2013，Stanier2013）。

\textbf{原文证据与缺口：}Birn2001针对已启动单片；L241图1--10中Hall与PIC的岛合并速率、压力阻力和尺寸依赖在弱导场下显著不同，强导场时某些率接近但拓扑和层宽仍不同；Stanier2013中磁通堆积、反弹及岛是否排出受装置几何控制。
\textbf{本报告综合判断：}局地快通量处理能力不足以唯一给出整个磁岛系统完成合并的时间。需要相同初边值、统一上游归一化和累计磁通诊断；这些论文没有共同建立跨系统的普适完成时间定律。

\subsection{有限三维X线的受抑长度是否普适，Hall模型能否独立复现}
\textbf{对应文章：}\citet{F659}（2019，F659，
\emph{Three-Dimensional Magnetic Reconnection With a Spatially Confined X-Line Extent}）；
\citet{F792}（2020，F792，
\emph{Scaling of Magnetic Reconnection With a Limited X-Line Extent}）。

\textbf{原文缺口：}F659图2--9及F792图2--5给出活跃/受抑区与整体率下降；两篇都是PIC研究，电子非对角压力仍支撑破冻结。约$10d_i$受抑尺度来自特定反平行序列，拟合参数、短层瞬态、导场及$\beta$推广仍有限制。
\textbf{本报告综合判断：}尚不能把该尺度写成标准Hall MHD的普适定律。有限长度与压力闭合是否相互改变，是由这两篇及F611共同引出的研究设计，并非三篇已经证实的结论。\citep{F611}

\subsection{离子加速怎样变成不可逆热化，能量怎样分给两物种}
\textbf{对应文章：}\citet{Yoo2013}（2013，Yoo2013，
\emph{Observation of Ion Acceleration and Heating during Collisionless Magnetic Reconnection in a Laboratory Plasma}）；
\citet{L272}（2021，L272）；
\citet{E512}（2011，E512）。

\textbf{原文缺口：}Yoo2013图2--4确立电势、加速与下游升温的空间联系，却保留碰撞、再磁化、波粒及边界等解释。L272图1--7提出电场梯度失磁化途径，但给定场能量展宽和PIC高温关联不能单独唯一确认所有条件下的热化机制；E512在强导场、低$\beta$排序下建立相混合与弱碰撞熵产链，范围又不同。新增核读的\citet{Drake2009Pickup}反平行拾取模型给中心矩温升标度，但观测系数偏低；其第16段与Yoon2023原文Discussion部分对电场净功和拾取主导程度的归因有分歧。\textbf{本报告判断：}需同轨道逐分量电功预算，不能以末速度、温度或Hall场图单独裁决。本轮修复的是这项比较在报告中缺失，并未宣称解决该研究分歧。
\textbf{11月12日专题补充：}双温Hall并非完全不能研究加热。\citet{L223}（2016）图3已给出两物种不同热点，\citet{Kaptanoglu2020}以三维双温Hall比较HIT-SI实验，\citet{TenBarge2019MMS}（2019）图1已输出两物种热能预算；但分别受超电阻/输运闭合、未取得的勘误及压力热流闭合限制。\citet{Tanabe2015MAST}的热点位置为前一类模拟提供实验约束，尚未唯一验证局地热化机制。

\textbf{比例必须区分：}\citet{F674}正式版图7的约2:1是离子/电子电场沉积率；同图热能变化约1:1，且含储存与边界焓流。\citet{Yoo2017Heating}非对称MRX沉积更接近；\citet{Yoon2023PhaseSpace}仍明确留下精确离子能量份额和有序/随机增能占比。\textbf{本报告综合判断：}不能用Hall总功为零否定物种增能，也不能用局地正电功替代不可逆热化。条件性两温预算已经可算；普适不可逆加热比仍需动理学信息或经独立检验的闭合，并分别排查模型缺失、数值误差、输运边界和诊断不足。独立专题源与PDF位于项目根目录\texttt{specific problems/离子加热问题/}。

\subsection{自然触发阈值与三维后期状态能否由同一模型预测}
\textbf{对应文章：}\citet{Birn2001}（2001，Birn2001）；
\citet{Cassak2005}（2005，Cassak2005）；
\citet{Cassak2010}（2010，Cassak2010）；
\citet{Shepherd2012}（2012，Shepherd2012）。

\textbf{原文缺口：}GEM基准跳过线性触发并明确保留三维问题；双稳态研究的触发与耗散、片层历史相关，简化模型未闭合层长和变化的上游场；Shepherd2012只在斜模明显相互作用前拟合传播，后期不再满足孤立X线图像。
\textbf{本报告综合判断：}这些研究各解决一个阶段，尚不能直接拼成从任意初始电流片到自然启动、三维扩展及饱和的无参数预测。第3章的启动线与第三维线应联合检验，而不能把阶段性公式无条件串接。
